% Patent Application LaTeX Template for Individual Inventors in Theoretical Domains
% Optimized for mathematical frameworks and abstract-concept inventions
% Incorporates special guidance for mechanism-based claim strategy and AI prompting

\documentclass[12pt,letterpaper]{article}

% Essential packages for patent formatting
\usepackage[top=1in, bottom=1in, left=1in, right=1in]{geometry} % Proper margins
\usepackage{setspace} % For 1.5 or double line spacing
\usepackage{times} % Times Roman font (or use \usepackage{helvet} for Arial)
\usepackage{enumitem} % For customizing lists/enumerations in claims
\usepackage{fancyhdr} % For headers
\usepackage{titlesec} % For section formatting
\usepackage{graphicx} % For including drawings
\usepackage{amsmath} % For mathematical equations
\usepackage{amsthm} % For theorems and proofs
\usepackage{algorithm} % For algorithm presentations
\usepackage{algpseudocode} % For pseudocode
\usepackage{listings} % For code listings

% Set the line spacing to 1.5 (or 2.0 for double spacing)
\onehalfspacing

% Configure headers and footers
\pagestyle{fancy}
\fancyhf{} % Clear all headers and footers
\renewcommand{\headrulewidth}{0pt} % Remove header rule
\fancyfoot[C]{\thepage} % Page number at bottom center

% Configure section formatting
\titleformat{\section}
  {\normalfont\bfseries\uppercase}
  {}
  {0em}
  {}
\titlespacing{\section}{0pt}{12pt}{6pt}

% Configure subsection formatting
\titleformat{\subsection}
  {\normalfont\bfseries}
  {}
  {0em}
  {}
\titlespacing{\subsection}{0pt}{12pt}{6pt}

% Custom paragraph format with additional spacing
\setlength{\parindent}{0pt}
\setlength{\parskip}{6pt}

% Custom enumeration for claims
\newlist{patentclaims}{enumerate}{3}
\setlist[patentclaims,1]{label=\arabic*., align=left, leftmargin=2em, itemsep=1em}
\setlist[patentclaims,2]{label=\arabic*., align=left, leftmargin=2em, itemsep=0.5em}
\setlist[patentclaims,3]{label=\arabic*., align=left, leftmargin=2em, itemsep=0.5em}

% Define theorem environment for mathematical proofs
\newtheorem{theorem}{Theorem}[section]
\newtheorem{definition}{Definition}[section]

\begin{document}

% ===== AI PROMPTING GUIDANCE =====
% This template contains special instructions for AI assistance in drafting patents
% for theoretical/mathematical domains. Look for the AI-PROMPT sections throughout
% the document for specific guidance on what to ask the AI to help with.

% ===== TITLE PAGE =====
\begin{center}
\Large\bfseries
[TITLE]
\end{center}

\vspace{0.5in}

% AI-PROMPT: Title Generation
% Generate a concise, technically precise title that clearly describes the core functional mechanism 
% of my theoretical invention. Focus on the specific computational or mathematical approach rather
% than results or benefits. The title should be 5-15 words and avoid terms like "new," "improved,"
% or other relative qualifiers. For mathematical inventions, emphasize the method or system rather than
% the underlying mathematical concept itself.
% Example good title: "Transfinite Ordinal Computing System With Hierarchical Shell Mapping"
% Example bad title: "Improved Method for Mathematical Calculations"

\vspace{1in}

\section*{CROSS-REFERENCE TO RELATED APPLICATIONS}

[CROSS-REFERENCE]

% AI-PROMPT: Cross-Reference
% If this is my first patent in this family, suggest appropriate text stating "Not Applicable."
% If I have related applications, help me properly reference them with their application numbers,
% filing dates, and relationship to this application. Use precise legal terminology for the
% relationship (continuation, continuation-in-part, divisional, etc.).

\section*{FIELD OF THE INVENTION}

[FIELD]

% AI-PROMPT: Field Definition
% Based on my description of the invention, craft a precise technical definition of the field 
% to which my theoretical invention pertains. Focus on the specific computational or mathematical
% domain rather than potential applications. For mathematical frameworks, specify the branch of
% mathematics and the particular problem-solving approach. Limit to 2-3 sentences that precisely
% position the invention within its technical field.
% Example: "The present invention relates to computational representations of transfinite ordinals, 
% and more particularly to systems and methods for implementing mathematical operations across 
% ordinal domains. Specifically, the invention pertains to data structures and algorithms for 
% bridging discrete and continuous mathematical spaces through ordinal-indexed computational structures."

\section*{BACKGROUND}

% AI-PROMPT: Background Construction
% Create a background section that establishes the technical foundation for my theoretical invention
% without overstating problems with prior approaches. For mathematical inventions, focus on:
% 1. The existing mathematical/computational approaches in this specific domain
% 2. The technical limitations of those approaches (describe objectively without disparagement)
% 3. The specific computational challenges these limitations create
% 4. Previous attempts to address these challenges and their technical trade-offs
% 
% Use precise technical language and avoid making statements about what is "impossible" or 
% "cannot be done" with prior approaches. Instead, describe specific computational inefficiencies
% or limitations. For theoretical frameworks, emphasize mathematical or computational constraints
% rather than application-specific limitations.
%
% End with a statement identifying the specific technical need that remains unaddressed,
% focusing on the mechanism rather than results.

\section*{SUMMARY}

% AI-PROMPT: Summary Construction
% Craft a summary that focuses on the core mechanism of my theoretical invention rather than its
% results or applications. For mathematical or computational inventions:
% 1. Begin with a statement that precisely articulates the invention in terms of its mechanism, 
%    not its purpose or benefits
% 2. Focus on "what the invention is" rather than "what the invention does"
% 3. Describe the central mathematical or computational approach without limiting language
% 4. Avoid focusing on specific results, instead emphasizing the underlying mechanism
% 5. If providing examples, use "in one embodiment" language to avoid limiting scope
%
% For theoretical frameworks, state the invention in terms of its core computational structure 
% or mathematical foundation. Align this summary with the broadest independent claim.
%
% Example structure: "The present invention provides a computational system for representing and 
% manipulating transfinite ordinals through a hierarchical shell data structure. In one aspect, 
% the invention comprises a specialized memory architecture for storing ordinal-indexed data with 
% distinct handling for finite and transfinite domains. The system implements bidirectional 
% mapping functions between ordinal indices and numerical approximations to enable efficient 
% computation across ordinal boundaries."

\section*{BRIEF DESCRIPTION OF THE DRAWINGS}

% AI-PROMPT: Drawing Description
% Create a comprehensive list of drawing descriptions that effectively represent my abstract 
% theoretical invention through visualizations. For mathematical frameworks, include:
% 1. Block diagrams showing system architecture or data flow
% 2. Flowcharts illustrating algorithmic processes
% 3. Mathematical relationship diagrams
% 4. Data structure representations
% 5. Graph or chart visualizations of theoretical relationships
% 6. Visual representations of abstract concepts
%
% For each figure, provide a clear description following the format "FIG. X is a [type] showing/illustrating
% [specific content]." Ensure all key aspects of the theoretical framework have visual representation,
% particularly the mechanism aspects that will be claimed.

\section*{DETAILED DESCRIPTION}

% AI-PROMPT: Detailed Description Construction
% Create a comprehensive detailed description of my theoretical invention that satisfies the enablement
% and written description requirements while maintaining maximum claim scope. For mathematical/computational
% inventions:
%
% 1. Begin with a general overview of the invention's mechanism without unnecessary limitations
%
% 2. Describe multiple embodiments showing different implementations of the core mechanism, 
%    demonstrating breadth while maintaining enablement
%
% 3. For mathematical concepts that might be considered abstract ideas:
%    - Focus on specific technical implementations and tangible structures
%    - Describe how the mathematical concepts integrate with concrete computing elements
%    - Emphasize practical applications and technical improvements
%    - Explicitly state technical advantages and improvements to computing technology
%
% 4. Include detailed mathematical derivations and theoretical foundations to satisfy the enablement
%    requirement, using theorems and proofs where appropriate
%
% 5. Provide pseudocode or algorithm descriptions for key computational processes
%
% 6. Show multiple specific applications of the theoretical framework to demonstrate practical utility
%
% 7. Include prophetic examples using present tense (not past tense) that demonstrate how the 
%    invention would be implemented in various contexts
%
% 8. For each mathematical concept or abstract idea:
%    - Describe specific hardware, memory, or computational structures involved
%    - Explain how the mathematical concept is specially adapted to a particular technological environment
%    - Demonstrate that the invention improves the functioning of a computer or other technology
%
% 9. For aspects potentially vulnerable to Section 101 challenges:
%    - Describe specific technical problems solved by the invention
%    - Articulate how the invention represents an improvement to computer technology
%    - Explain why the solution is necessarily rooted in computer technology
%    - Detail technical advantages beyond conventional approaches
%
% 10. Provide fallback positions with increasingly narrower embodiments to support dependent claims

\section*{MATHEMATICAL FOUNDATIONS}

% AI-PROMPT: Mathematical Foundations
% Create a detailed section establishing the mathematical foundations of my theoretical invention.
% This section should:
% 
% 1. Define all mathematical terms, variables, and symbols used in the invention
% 
% 2. Present key mathematical theorems and their proofs that underpin the invention
% 
% 3. Demonstrate how these mathematical concepts translate to computational implementations
%
% 4. Show derivations that connect the mathematical theory to practical applications
%
% 5. Provide mathematical proofs of key properties or behaviors of the invention
%
% 6. Define boundary conditions and constraints on the mathematical framework
%
% 7. Present formal definitions using precise mathematical notation
%
% For each mathematical concept, include both the theoretical foundation and at least one 
% concrete implementation example showing how the mathematics translates to a computational system.
% This section should be rigorous enough to satisfy enablement requirements while establishing
% the theoretical validity of the invention.

\section*{PRACTICAL APPLICATIONS}

% AI-PROMPT: Practical Applications
% Create a section detailing specific practical applications of my theoretical invention that
% demonstrate its patent eligibility under Section 101. For each application:
%
% 1. Describe a specific technical problem that the invention addresses
%
% 2. Explain how the invention solves this problem through its unique mechanism
%
% 3. Demonstrate why this application represents a technological improvement
%
% 4. Explain why the solution is necessarily rooted in computer technology
%
% 5. Provide specific examples of how the invention improves the functioning of a computer
%    or other technology in this application domain
%
% Ensure each application clearly shows the invention is not merely an abstract idea but
% is integrated into a practical application with specific technological improvements.
% Focus on computational efficiency, improved accuracy, reduced resource usage, or other
% technical benefits rather than business or user experience improvements.

\section*{IMPLEMENTATION ARCHITECTURES}

% AI-PROMPT: Implementation Architectures
% Create a detailed section describing specific technical implementations of my theoretical invention
% across different hardware and software architectures. For each architecture, include:
%
% 1. Specific hardware components and their configurations
%
% 2. Memory structures and data organization
%
% 3. Processing flow and computational mechanisms
%
% 4. Interfaces and communication protocols
%
% 5. Specific technical advantages of the implementation
%
% Include implementations for:
% - General-purpose computing systems
% - Specialized hardware accelerators
% - Cloud computing environments
% - Edge computing devices
% - High-performance computing clusters
%
% For each architecture, explain how the mathematical or theoretical aspects of the invention
% are specifically implemented in that environment, providing sufficient technical detail
% to satisfy the enablement requirement.

\section*{EXEMPLARY EMBODIMENTS}

% AI-PROMPT: Exemplary Embodiments
% Create a section with multiple detailed embodiments of my theoretical invention in different
% contexts. For each embodiment:
%
% 1. Describe a specific implementation scenario (use present tense for prophetic examples)
%
% 2. Detail the exact mechanism of operation in this specific context
%
% 3. Provide technical parameters, configurations, or settings
%
% 4. Explain specific advantages in this implementation
%
% 5. Include pseudocode, algorithms, or processing steps
%
% Provide a range of embodiments from broad to narrow to support different claim scopes,
% ensuring the broadest embodiments remain enabled while narrower embodiments provide
% fallback positions. Include specific numerical ranges and parameters where appropriate,
% but avoid unnecessarily limiting language except in the narrowest embodiments.

\newpage
\section*{CLAIMS}

% AI-PROMPT: Claims Construction
% Construct a hierarchical claim set for my theoretical invention that maximizes protection while
% ensuring patent eligibility. Follow these specific approaches for theoretical/mathematical inventions:
%
% 1. Structure independent claims around mechanisms rather than results, focusing on:
%    - Specific data structures or computational architectures
%    - Concrete processing steps or transformations
%    - Tangible implementations of mathematical concepts
%
% 2. Ensure abstract ideas are integrated with practical applications by:
%    - Tying mathematical concepts to specific technological improvements
%    - Including elements that show the invention improves computer functionality
%    - Demonstrating the solution is necessarily rooted in computer technology
%
% 3. Create a hierarchical claim structure:
%    - Broadest claims: Focus on core mechanisms without application-specific limitations
%    - Mid-level claims: Add specific computational advantages or implementations
%    - Narrow claims: Include specific applications or optimizations
%
% 4. For method claims:
%    - Begin with concrete initialization or configuration steps
%    - Include specific data transformation or processing operations
%    - End with generating or outputting tangible results or structures
%
% 5. For system claims:
%    - Define specific memory structures or data organizations
%    - Include processing components with specific functions
%    - Establish clear relationships between components
%
% 6. Avoid claims focusing solely on mathematical algorithms by:
%    - Integrating them with specific technical elements
%    - Showing how they solve a specific technical problem
%    - Demonstrating they improve computer functionality
%
% 7. Include both method and system claims to provide multiple protection angles

\begin{patentclaims}
    \item A system for computational representation of transfinite ordinals, comprising:
    \begin{itemize}[leftmargin=2em]
        \item a memory architecture configured to store ordinal-indexed data with separate storage mechanisms for finite and transfinite ordinal indices;
        \item a processing unit configured to perform operations on the ordinal-indexed data; and
        \item an ordinal mapping module configured to convert between symbolic ordinal representations and numerical approximations,
        \item wherein the system implements bidirectional data flow between ordinal domains to enable mathematical operations across ordinal boundaries.
    \end{itemize}
    
    \item The system of claim 1, wherein the memory architecture implements a hierarchical structure with exponentially increasing capacity for higher ordinal indices.
    
    \item The system of claim 2, wherein the hierarchical structure comprises:
    \begin{itemize}[leftmargin=2em]
        \item a first memory region storing data associated with finite ordinals using direct addressing;
        \item a second memory region storing data associated with countable transfinite ordinals using indexed mapping; and
        \item a third memory region storing data associated with uncountable transfinite ordinals using hash-based addressing.
    \end{itemize}
    
    \item The system of claim 1, wherein the processing unit implements ordinal arithmetic operations comprising:
    \begin{itemize}[leftmargin=2em]
        \item ordinal addition that preserves non-commutativity properties;
        \item ordinal multiplication with proper left and right distributivity; and
        \item ordinal exponentiation for finite and transfinite exponents.
    \end{itemize}
    
    \item A method for computational representation and manipulation of transfinite ordinals, comprising:
    \begin{itemize}[leftmargin=2em]
        \item initializing a memory structure with distinct regions for storing finite and transfinite ordinal-indexed data;
        \item implementing bidirectional mapping functions between symbolic ordinal representations and numerical approximations;
        \item performing ordinal arithmetic operations using the bidirectional mapping functions; and
        \item generating computational results that preserve ordinal arithmetic properties.
    \end{itemize}
    
    \item The method of claim 5, further comprising:
    \begin{itemize}[leftmargin=2em]
        \item detecting mathematical operations that cross ordinal boundaries;
        \item automatically selecting appropriate computational representations based on ordinal domains; and
        \item applying domain-specific optimization for operations within each ordinal domain.
    \end{itemize}
    
    \item The method of claim 5, wherein performing ordinal arithmetic operations comprises:
    \begin{itemize}[leftmargin=2em]
        \item determining ordinal relationship types between operands;
        \item selecting specialized algorithms based on ordinal relationship types; and
        \item implementing ordinal arithmetic using the selected algorithms.
    \end{itemize}
    
    \item A non-transitory computer-readable medium storing instructions that, when executed by a processor, cause the processor to:
    \begin{itemize}[leftmargin=2em]
        \item initialize a hierarchical memory structure for representing ordinal-indexed data;
        \item implement bidirectional mapping functions between ordinal indices and numerical values;
        \item perform mathematical operations across ordinal domains using the bidirectional mapping functions; and
        \item generate computational results that maintain mathematical consistency across ordinal boundaries.
    \end{itemize}
    
    \item The non-transitory computer-readable medium of claim 8, wherein the instructions further cause the processor to:
    \begin{itemize}[leftmargin=2em]
        \item implement an ordinal-based regularization scheme for mathematical operations;
        \item automatically select computational precision based on ordinal domain; and
        \item dynamically allocate computational resources based on ordinal complexity.
    \end{itemize}
    
    \item The non-transitory computer-readable medium of claim 8, wherein the mathematical operations comprise at least one of:
    \begin{itemize}[leftmargin=2em]
        \item transfinite summation across ordinal domains;
        \item ordinal-indexed matrix operations;
        \item shell coupling calculations between ordinal-indexed structures; and
        \item bidirectional flow computations between nested ordinal domains.
    \end{itemize}
\end{patentclaims}

\newpage
\section*{ABSTRACT}

[ABSTRACT]

% AI-PROMPT: Abstract Creation
% Create a concise, technically-focused abstract for my theoretical invention that highlights
% the mechanism rather than results or benefits. For mathematical/computational inventions:
%
% 1. Begin with a sentence clearly stating what the invention is in terms of its mechanism
%
% 2. Briefly describe the core technical approach or computational architecture
%
% 3. Include 1-2 sentences on how the mechanism works, focusing on the "how" not the "why"
%
% 4. Briefly mention a key technical advantage that relates to computing technology
%
% 5. Stay under 150 words total and avoid legal terminology
%
% For theoretical frameworks, focus on the computational implementation rather than the 
% underlying mathematical concepts. Emphasize technical elements that demonstrate the 
% invention is more than an abstract idea.
%
% Example: "A computational system for representing and manipulating transfinite ordinals 
% through hierarchical memory structures with distinct handling for finite and transfinite 
% domains. The system implements bidirectional mapping functions that convert between symbolic 
% ordinal representations and numerical approximations, enabling mathematical operations 
% across ordinal boundaries. These operations maintain mathematical consistency while achieving 
% computational tractability through specialized algorithms for different ordinal relationship 
% types. The hierarchical memory architecture provides exponentially increasing capacity for 
% higher ordinal indices, enabling efficient storage and retrieval of ordinal-indexed data 
% across both finite and transfinite domains."

\newpage
\section*{DRAWINGS}

[DRAWINGS]

% AI-PROMPT: Technical Drawing Guidance
% Create technical drawings that effectively visualize my abstract theoretical invention through
% diagrams, flowcharts, and visual representations. For mathematical/computational frameworks:
%
% 1. System Architecture Diagram: Show the overall structure of the computational system with
%    components clearly labeled
%
% 2. Memory Structure Diagram: Illustrate the hierarchical memory organization showing how
%    different ordinal domains are represented
%
% 3. Data Flow Diagram: Demonstrate how data moves through the system during key operations
%
% 4. Algorithm Flowcharts: Present the logical flow of key algorithms with decision points
%    and processing steps
%
% 5. Mathematical Relationship Visualization: Show how mathematical concepts relate visually
%    through graphs, charts, or other visual representations
%
% 6. Ordinal Mapping Diagram: Illustrate how the system maps between symbolic and numerical
%    representations across ordinal boundaries
%
% 7. Implementation Architecture Diagrams: Show how the system would be implemented in different
%    computing environments
%
% For each drawing, use clear labels, consistent numbering, and visual organization that helps
% explain the abstract concepts in a concrete way. Focus on aspects that demonstrate the invention
% is tied to specific technological implementations rather than purely abstract ideas.

% Example placeholder for drawings
\begin{figure}[h]
\centering
\framebox{FIG. 1: System Architecture Diagram showing the overall structure of the ordinal computation system}
\caption*{FIG. 1}
\end{figure}

\begin{figure}[h]
\centering
\framebox{FIG. 2: Hierarchical Memory Structure for Ordinal-Indexed Data}
\caption*{FIG. 2}
\end{figure}

\begin{figure}[h]
\centering
\framebox{FIG. 3: Bidirectional Mapping Function Between Symbolic and Numerical Representations}
\caption*{FIG. 3}
\end{figure}

\begin{figure}[h]
\centering
\framebox{FIG. 4: Flowchart of Ordinal Arithmetic Operations Across Domain Boundaries}
\caption*{FIG. 4}
\end{figure}

\begin{figure}[h]
\centering
\framebox{FIG. 5: Visual Representation of Ordinal Relationships and Coupling}
\caption*{FIG. 5}
\end{figure}

\begin{figure}[h]
\centering
\framebox{FIG. 6: Implementation Architecture for High-Performance Computing Environment}
\caption*{FIG. 6}
\end{figure}

\end{document}